% !TeX root = ../tesis.tex
%%%%%%%%%%%%%%%%%%%%%%%%%%%%%%%%%%%%%%%%%%%%%%%%%%%%%%%%%%%%%%%%%%%
% ANEXO G — Recursos Visuales: Escenarios con el Anillo Periférico
%           Interior (API-I)
% Issue #25. Paralelo al Anexo F (red actual).
%
% Regla del autor (2026-10-18): aquí solo entran los recursos de las
% secciones ya redactadas del Cap. 6; los demás se agregan al
% terminar cada sección o subsección. Mapas a página completa
% (girados 90°, [p]).
%
% Estado:
%   G.1  Cartografía QGIS — Propuesta del API-I
%        · Propuesta_Base_Anillar_Metro (§6.1.2, fig:propuesta-anillo)
%   Pendientes (al redactar §6.3–§6.4):
%        · Propuesta_Base_Anillar_MB, Propuesta_Anillar_Extendido_MB,
%          Propuesta_Anillar_Extendido_Metro (§6.3.1)
%        · Clasificación Garibelt por escenario (QGIS)
%        · Dashboards por escenario y comparativos (Tableau)
%%%%%%%%%%%%%%%%%%%%%%%%%%%%%%%%%%%%%%%%%%%%%%%%%%%%%%%%%%%%%%%%%%%

Este anexo reúne los recursos visuales de los dos escenarios del
\textsc{api-i}: el Caso~2 (Anillo MB) y el Caso~3 (Anillo Metro). Cumple para
el Capítulo~\ref{ch:analisis-red-anillo} la función que el Anexo~F cumple para
la red actual: guarda a página completa los mapas que el capítulo presenta en
tamaño reducido. Los mapas se produjeron con \textbf{QGIS~3.x} en la
proyección \textsc{epsg}:32614 (UTM zona~14N), con los datos de la corrida
2026-10-04.

% =============================================================================
\section{Cartografía \textsc{QGIS} — Propuesta del \textsc{api-i}}
\label{anx:g-qgis-propuesta}
% =============================================================================

% --- G.1.1 -------------------------------------------------------------------
\subsection{Trazo del \textsc{api-i}}
\label{anx:g-propuesta-base}

\begin{figure}[p]
  \centering
  \includegraphics[angle=90,height=0.82\textheight]{%
    Figures/Mapas/Qgis_Resources/RedConAnillo/Escenario_Metro/Base/Propuesta_Base_Anillar_Metro}
  \caption[Trazo del \textsc{api-i} sobre la red de la \zmvm{} (página
  completa)]{Trazo del \textsc{api-i} sobre la red de transporte de la
  \zmvm{}. Los Casos~2 (Anillo MB) y~3 (Anillo Metro) comparten el trazo y
  las 94 estaciones. Versión reducida en la
  Figura~\ref{fig:propuesta-anillo}. Corrida: 2026-10-04.}
  \label{fig:anx-g-propuesta-base}
  \fuentefigura{Elaboración propia con \texttt{Transport-gis-zmvm-mjg};
    datos: \texttt{Apímetro} y \texttt{VFTModel}~\citep{cabrera2025apimetro, cabrera2026VFTModel}.}
\end{figure}
